\documentclass[10pt,a4paper]{amsart}
\usepackage[margin=19mm]{geometry}
\usepackage{amsmath,amssymb,mathtools}
\usepackage[scr=rsfs,cal=euler]{mathalfa}
\usepackage{xfrac}
\usepackage[unicode,hidelinks,bookmarks=false]{hyperref}
\newcommand{\ie}{i.e.\ }
\newcommand{\cf}{cf.\ }
\newcommand{\resp}{respectively\ }
\newcommand{\Subsection}{\subsection}
\providecommand{\nobreakdash}{\nobreak\hbox{-}}
\setlength{\emergencystretch}{2em}
\multlinegap0pt
\usepackage{amssymb}
\usepackage{mathtools}
\usepackage[shortlabels]{enumitem}
\usepackage{xcolor}
\usepackage{tikz}
\newtheorem{theorem}{Theorem}
\newtheorem{lemma}{Lemma}
\newcommand{\Aut}{\operatorname{Aut}}
\newcommand{\R}{\mathbb R}
\newcommand{\Z}{\mathbb Z}
\newcommand{\cX}{\mathcal X}
\newcommand{\elltwo}{\ell^2}
\title{Regular hyperbolic tilings have no $\ell^2$ eigenfunctions}
\author{Asaf Nachmias}
\date{Source: arXiv:2609.11857v1 (CC BY 4.0)}
\begin{document}
\maketitle

\noindent\emph{Source and licence.} Regular hyperbolic tilings have no $\ell^2$ eigenfunctions. Version arXiv:2609.11857v1 (CC BY 4.0), distributed by arXiv under the Creative Commons Attribution 4.0 International licence, http://creativecommons.org/licenses/by/4.0/, as recorded in the arXiv licence metadata for this exact version and shown on the abstract page https://arxiv.org/abs/2609.11857v1. Full licence notice: \texttt{licenses/arxiv-2609.11857-CC-BY-4.0.txt}.\par\smallskip

\noindent\textit{Editorial scope.} Counted unit: the article's theorem thm:criterion (source0.tex lines 248--254). The article's complete original proof of it is delivered with it (source0.tex lines 453--530, one \texttt{proof} environment of 78 lines, which the article places away from the statement). No mathematics is written, repaired or re-derived: the statement and the proof are the article's own lines, in the article's own notation, and no retained line is substituted or re-flowed.  Retained uncounted as the source-local context the proof needs: lines 381--394; lines 399--416.  Nothing was omitted from any retained span, so the package carries no omission record.  No retained line contains a citation, so the wrapper carries no bibliography and no bibliography entry.  No retained line contains a figure environment, an image-inclusion command or drawing code, so no image is delivered and the package declares no asset dependency.  Editorial formatting only, in addition: an AMS \texttt{amsart} preamble that restates the article's own declarations and macros used by the retained lines, namely the environments \texttt{theorem}, \texttt{lemma} on separate counters, together with the standard packages those lines need.  The retained environments print in this document as Theorem 1, Lemma 1 in order of appearance, whereas the article prints the same items with the numbering of its own full document.  The complete native source of the article as distributed in the arXiv e-print for this exact version is bundled unchanged as \texttt{sources/210057/source0.tex}.
\begin{theorem}\label{thm:Bordenave} 
Let $(G,o)$ be a unimodular rooted graph that is almost surely bounded degree, and let $\eta$ be an invariant labeling having $L$ distinct values. Then for any $\lambda\in \R$
\[
\mathbb E\,\mu_o(\{\lambda\})
\le
\mathbb P(o\text{ is level or bad}) \, . \]
%+
%\sum_j
% \mathbb E\!\left[
%   \mu^{(j)}_o(I)\,1_{\{o\text{ lies in a level block of label }j\}}
% \right] \, .

%where $\mu^{(j)}_o$ is the spectral measure for the restriction of the operator to the level-$j$ vertices.
\end{theorem}

\begin{lemma}\label{lem:height} Let $H\leq \Aut(\cX_{p,q})$ be a subgroup which acts freely and cocompactly on the vertices. If $u:E(H\backslash \cX)\to \Z$ is an $H$-cocycle, then there exists a function
\[
h:V(\mathcal X) \longrightarrow\Z
\]
and a homomorphism
\[
\phi:H\longrightarrow\Z
\]
such that for every oriented edge $(w,v)\in E(\mathcal X)$
\[
h(v)-h(w)=\tilde{u}(w,v)
\]
and for every $g\in H$ and $v\in V(\cX)$ 
\[
h(gv)-h(v)=\phi(g) \, .
\]
%The image of $\phi$ is the subgroup of $\Z$ generated by all the values of $u$.  In particular, $u$ is primitive if and only if $\phi$ is onto $\Z$.
\end{lemma}

\begin{theorem}\label{thm:criterion}
Let $\mathcal X_{p,q}$ be the $(p,q)$ regular tiling and let $H\leq\Aut(\mathcal X_{p,q})$ act freely and cocompactly on the vertices. Let $G$ be the $1$-skeleton of $\mathcal X_{p,q}$. If $\mathcal X_{p,q}$ admits a monic $H$-cocycle, then the adjacency operator of $G$ has no nonzero square-integrable eigenfunction, that is, for every $\lambda\in\R$
\[
\ker(A-\lambda I)=\{0\}
\qquad\text{in }\elltwo(G) \, .
\]
\end{theorem}

\begin{proof}[Proof of Theorem \ref{thm:criterion}]
Let $\mathcal X$ denote the simply connected complex and $G$ its
$1$-skeleton. Write $Q=H\backslash\mathcal X$, number its vertices
$0,\ldots,d-1$, and let $\tau:V(X)\to \{0,\ldots,d-1\}$ record the quotient vertex of each vertex of $G$. Let $u$ be the monic $H$-cocycle,  $\widetilde u$ be its lift to $G$, and set
$k=\max u$. Since $u$ is nonzero and antisymmetric, $k\geq 1$ and
$|\widetilde u(e)|\leq k$ for every oriented edge $e$.

Recall that $T_u$ has an edge $L_iR_j$ for each oriented quotient edge
from $i$ to $j$ whose $u$-value is $k$, with parallel edges retained. Since $T_u$ has a perfect matching, there exists a permutation $\pi$ of $\{0,\ldots, d-1\}$ so that $L_i R_{\pi(i)}$ is the matching. Since this matching is unique, we can relabel the set $\{0,\ldots, d-1\}$ so that if $L_i R_j$ is a non matching edge of $T_u$, then $j < \pi(i)$. Indeed, this follows since the directed graph on the set $\{0,\ldots, d-1\}$ obtained by adding the edge $j \to \pi(i)$ when $L_iR_j$ is a non matching edge of $T_u$, is acyclic due to the uniqueness of the perfect matching (otherwise we could replace edges of the matching along that cycle and obtain a contradiction to uniqueness).

Apply Lemma \ref{lem:height} and let $h:V(G)\to \Z$ and $\phi:H\to\Z$ denote the corresponding equivariant height functions and homomorphism, respectively.
% Let $M$ be its unique perfect matching. There is a permutation
% $\pi$ of $\{1,\ldots,d\}$ such that the matching edge incident to
% $L_i$ ends at $R_{\pi(i)}$. Write $e_i$ for the corresponding oriented
% quotient edge.
% For each non matching edge $L_i R_j$ in $T_u$ draw a directed edge $j \to \pi(i)$ on the set $\{1,\ldots, d\}$. The resulting directed graph is acyclic. Indeed, if there
% was a cycle, we can replace the matching edges along that cycle by non-matching edges and produce
% a second perfect matching. Therefore, we may assign distinct integers $r_1,\ldots, r_d \in \{0,\ldots, d-1\}$ to the vertices so that $r_j \leq r_{\pi(i)}$ for every non matching edge $L_i R_j$ of $T_u$.
Define the modified height function $\xi:V(G)\to\Z$ by
\begin{equation*}\label{eq:repaired-modified-height}
  \xi(v)=d\,h(v)+\tau(v).
\end{equation*}
Since $\tau(gv)=\tau(v)$ for any $g\in H$, we have
\begin{equation}\label{eq:repaired-modified-equivariance}
  \xi(gv)=\xi(v)+d\,\phi(g),\qquad g\in H.
\end{equation}

We now observe that every vertex is a prodigy for $\xi$. If $L_i R_{\pi(i)}$ is a matching edge of $T_u$, then it corresponds to a quotient edge $(w,v)$ with $u$-value $k$ so that $\tau(w)=i$ and $\tau(v)=\pi(i)$. By Lemma \ref{lem:height} we have $h(v)-h(w)=k$ and
$$ \xi(v) - \xi(w) = dk + \pi(i) - i \geq dk - (d-1) \geq 1 \, .$$
Let $z \neq v$ be any other neighbor of $w$. If $h(z) - h(w) < k$, since $h$ is $\Z$-valued we get $h(v)-h(z) \geq 1$, hence
$$ \xi(v) - \xi(z) = d(h(v)-h(z)) + \tau(v) - \tau(z) \geq d - (d-1) = 1\, .$$
Otherwise $h(z)-h(w)=k$ and so the edge $(w,z)$ is maximal and present in $T_u$. Since $H$ acts freely on the vertices, the edges $(w,z)$ and $(w,v)$ are different quotient edges. Therefore $\tau(z) < \tau(v)=\pi(i)$. Hence
$$ \xi(v) - \xi(z) = \pi(i) - \tau(z) > 0 \, ,$$
so $v$ is prodigy for $\xi$. 

Next, fix an integer $L>2d(k+1)$, let $\omega$ be drawn uniformly from
$\{0,\ldots,L-1\}$ and define
$$ \eta_\omega(v)=(\xi(v)+\omega)\bmod L \, . $$
The constant $d(k+1)$ is chosen since $|\xi(x)-\xi(y)|\leq d(k+1)$ whenever $x \sim y$. By~\eqref{eq:repaired-modified-equivariance}, for each $g\in H$,
\[
  \eta_\omega(gv)
  =\bigl(\xi(v)+\omega+d\,\phi(g)\bigr)\bmod L.
\]
Adding a fixed residue preserves the uniform law of $\omega$. Hence equivariance guarantees that the law of the entire labeling is $H$-invariant. Furthermore, it is standard to verify unimodularity. Indeed, choose  representatives $v_1,\ldots,v_d$ of the vertex
 orbits, and independently of $\omega$, choose $o$ uniformly from
 these representatives. We verify that $(G,o,\eta_\omega)$ is
 unimodular. Let $F$
 be any nonnegative measurable mass transport, invariant under
 isomorphisms of doubly rooted marked graphs. Freeness implies that
 each vertex has a unique representation $gv_j$ with $g\in H$. Thus
 \begin{align*}
   \mathbb E\sum_{x\in V(G)}F(G, \eta_\omega, o,x)
  &=\frac1d\sum_{i,j=1}^d\sum_{g\in H}
    \mathbb E_\omega F(G, \eta_\omega,v_i,gv_j)\\
  &=\frac1d\sum_{i,j=1}^d\sum_{g\in H}
    \mathbb E_\omega F(G,\eta_\omega,g^{-1}v_i,v_j) =\mathbb E\sum_{x\in V(G)}F(G, \eta_\omega,x,o).
\end{align*}
The second equality uses $H$-invariance of the labeling law; the
last uses the bijection $g\mapsto g^{-1}$ and the orbit
decomposition. 
%All sums are justified by nonnegativity. This is
%the mass-transport identity.


% Also, putting $B=dk+d-1$, we have
% \begin{equation}\label{eq:repaired-edge-bound}
%   |\xi(x)-\xi(y)|\leq B\qquad\text{whenever }x\sim y.
% \end{equation}

Let $v$ and its witness $w$ be as above. Since $v$ is prodigy, for any vertex $y\neq v$ that is either $w$ or one of $w$'s neighbors we have that $1 \leq \xi(v)-\xi(y) \leq 2d(k+1)$. If $\eta_\omega(v) = j \geq 2d(k+1)$, then $\eta_\omega(y) = j - (\xi(v)-\xi(y))$ is in $\{0,\ldots, j-1\}$, so $v$ is still prodigy for $\eta_\omega$ with the same witness. It follows that 
$$ \mathbb P(o\text{ is level or bad})\leq\frac{2d(k+1)}{L} \, ,$$
from which Theorem \ref{thm:Bordenave} implies that for any $\lambda\in \R$
we have
$$ \mathbb E\mu_o(\{\lambda\})
  \leq\mathbb P(o\text{ is level or bad})
  \leq\frac{2d(k+1)}{L} \, ,$$
but since our graph is vertex transitive and deterministic, we get that $\mu_v(\{\lambda\}) \leq 2d(k+1)/L$ for any $v\in V$. Taking $L\to\infty$ gives that $\mu_v(\{\lambda\})=0$ for all vertices. We conclude that there are no eigenfunctions. 
\end{proof}

\end{document}
